%% sections/discussion.tex -- Phase 7.1. Register: CLAUDE.md S9.2. The framework (tab_framework)
%% is the centre; "validate locally" governs the per-axis priors. Two threads: (1) consistency
%% sacrificed (SD x length only moderate/heterogeneous -- neither pure-mechanical nor pure-intrinsic),
%% (2) the gold lesson. The SemEval unseen_domain observation enters LABELLED post-hoc, separate
%% from the confirmatory RQ results.

\section{Discussion: The Configuration Framework}
\label{sec:discussion}

% auto-generated by make_phase5_tables.py -- do not hand-edit
\begin{table*}[!tp]\centering\small
\caption{The configuration decision guide (Section~\ref{sec:framework}). Each row pairs a grading scenario with a recommended setting and the trade-off behind it.}\label{tab:framework}
\begin{tabular}{p{2.8cm}p{4.2cm}p{9.2cm}}
\hline
Scenario & Recommended & Trade-off in the rule \\
\hline
\raggedright Grading short answers & Reasoning off; the cheapest model & The reasoning gain is small and only sometimes significant (Qwen on Mohler and SemEval, not the other models), and reasoning off is 10--925$\times$ cheaper \\[3pt]
\hline
\raggedright Grading code, the baseline collapses to zero & Reasoning on, or switch to a discriminating model & Reasoning recovers the discrimination the collapsed baseline lost (Qwen $+0.113$ QWK on PT-CS-verified code), but at about 800$\times$ the tokens and worse run-to-run consistency, so switching model is often the better fix \\[3pt]
\hline
\raggedright Grading code, the baseline does not collapse & Reasoning off; on only if local validation justifies it & Without the collapse the reasoning gain is inconsistent and not predictable in advance, so off is the safe default, saving cost (about 94--192$\times$ the tokens) and keeping consistency \\[3pt]
\hline
\raggedright Consistency or fairness is critical & Reasoning off & Reasoning on is less reproducible, partly a length effect and partly reasoning itself \\[3pt]
\hline
\raggedright A rubric or reference answer is available & Supply it as guidance (rubric for code, reference answer for short answers) & Guidance was the one change that reliably raised agreement: it helped across datasets (RIAYN $+0.123$ QWK, PT-CS-verified code $+0.146$), at the cost of only the extra prompt tokens \\[3pt]
\hline
\raggedright The rubric has several criteria & Holistic scoring (one overall score) & On PT-CS-verified code, scoring criterion by criterion did not reliably change agreement ($-0.022$ QWK, not significant) at a higher token cost \\[3pt]
\hline
\raggedright The exam has several questions & Grade in blocks of a few questions, one call per block & On PT-CS-verified code, grouping questions into one call barely changed agreement (question-by-question to whole-exam $-0.060$ QWK, not significant), so batching them wins on cost by cutting the per-call overhead \\[3pt]
\hline
\raggedright Grading many answers in sequence & Fresh conversation per block; no history carried across the exam & A shared, accumulating conversation made grading stricter and answer order added noise (a 20-submission sub-study) \\[3pt]
\hline
\end{tabular}
\end{table*}


The study's deliverable is the decision guide in Table~\ref{tab:framework}, and its rows should be
read as priors, not guarantees. Each row maps a task characteristic to a recommended starting
configuration and states the trade-off inside the rule, with the supporting evidence in
Section~\ref{sec:results}. The bottom row governs the others. Transfer was at best partial (the top public
configuration was not refuted, but the verified samples below it could not separate the ranking from
noise), effects were model-conditional (the reasoning gain) and gold-conditional (the rubric benefit),
and the gains below the top were not predictable a priori. A deployment should therefore adopt the row that matches its task as a starting
point and validate it locally, against a reference it trusts, before relying on it.

The code rows are keyed on observable model behaviour rather than model identity, and that choice is
itself a finding. Whether reasoning helps a code grader depended not on the model's name but on
whether its baseline collapsed to zero. The collapse is the one case the data explain (reasoning is
the lever, or switch model); off the collapse, the gains were scattered across cells with no
pattern a practitioner could anticipate. Identity-keyed advice would also age with every model
release; behaviour-keyed advice survives it, because the diagnostic (run the cheap baseline, look at
the grade distribution) is a one-afternoon local check. This is the framework's answer to the
fast-moving model landscape, and it is why the unpredictability below the collapse reinforces, rather
than weakens, the governing rule.

The first cross-cutting thread is that reasoning sacrifices consistency, and the loss is not fully
explained by mechanics. A natural deflation would be that longer outputs simply give a
non-deterministic backend more tokens to diverge on, making the inconsistency a length artefact. The
data support that only partially: the correlation between an item's grade variability and its output
length is moderate and heterogeneous across cells (Spearman median 0.46, range 0.11 to 0.78, weakest
on Mohler at 0.11 to 0.28). Part of the
consistency loss is therefore mechanical, length interacting with backend non-determinism, and part
behaves like a property of reasoning itself; we claim neither pure story. For the grader's user the
distinction matters less than the fact: the same answer, graded twice under identical settings,
disagrees with itself more when the model reasons, which is why the consistency-critical row of
Table~\ref{tab:framework} recommends reasoning off.

The second thread is the gold lesson. The most consequential error in this study would have been one
we nearly published: concluding that rubrics do not help code grading, when the null was an artefact
of an unvalidated reference (Section~\ref{sec:gold}). The curation that prevented it cost nothing but
honesty: a mechanical, auditable criterion (does the final grade differ from the criterion sum?)
separated grades with evidence of human review from the rest. Institutions evaluating LLM graders on
their own historical grades should assume that reference quality is a first-order variable, audit it
before measuring anything, and prefer a smaller trusted reference over a larger doubtful one. For the
assessment platform this study feeds, the design consequence is two floors: grade provenance must be
traceable (who validated this grade?), and the grading standard itself must be calibrated, because
validation evidence and rigour are different properties.

One observation enters as exploratory, post hoc, and outside the pre-registered questions. With
reasoning on, SemEval discrimination improved most on the unseen-domain split (AUROC $+0.013$ to
$+0.039$ across all four models) while the seen and unseen-questions splits fell for the open
models. If it replicates, it suggests reasoning helps most where surface familiarity helps least,
which would refine the short-answer row of the framework; we flag it for confirmatory follow-up and
build nothing on it here.

For practice, the framework reduces to three habits. Start from the row that matches the task, not
from the newest model. Measure consistency and cost alongside agreement, because two of the three
are routinely omitted and both moved against reasoning here. And before deployment, re-run the cheap
baseline locally: the one-afternoon check (grade distribution, agreement on a trusted sample,
repetition variance) is what converts these priors into a defensible configuration choice.
